
\documentclass{beamer}

%\usepackage[table]{xcolor}
\mode<presentation> {
  \usetheme{Boadilla}

%  \usetheme{Pittsburgh}
%\usefonttheme[2]{sans}
\renewcommand{\familydefault}{cmss}
%\usepackage{lmodern}
%\usepackage[T1]{fontenc}
%\usepackage{palatino}
%\usepackage{cmbright}
 \useinnertheme{rectangles}
}

\setbeamercolor{normal text}{fg=black}
\setbeamercolor{structure}{fg= blue}
\definecolor{trial}{cmyk}{1,0,0, 0}
\definecolor{trial2}{cmyk}{0.00,0,1, 0}
\definecolor{darkgreen}{rgb}{0,.4, 0.1}
\usepackage{array}
\beamertemplatesolidbackgroundcolor{white}  \setbeamercolor{alerted
text}{fg=red}
\usepackage{tcolorbox}
\setbeamertemplate{caption}[numbered]\newcounter{mylastframe}


\usepackage{tikz}
\usetikzlibrary{arrows}
\usepackage{colortbl}

\renewcommand{\familydefault}{cmss}
%\usepackage[all]{xy}

\usepackage{tikz}
\usepackage{lipsum}

 \newenvironment{changemargin}[3]{%
 \begin{list}{}{%
 \setlength{\topsep}{0pt}%
 \setlength{\leftmargin}{#1}%
 \setlength{\rightmargin}{#2}%
 \setlength{\topmargin}{#3}%
 \setlength{\listparindent}{\parindent}%
 \setlength{\itemindent}{\parindent}%
 \setlength{\parsep}{\parskip}%
 }%
\item[]}{\end{list}}
\usetikzlibrary{arrows}

\usecolortheme{lily}

\newtheorem{com}{Comment}
\newtheorem{lem} {Lemma}
\newtheorem{prop}{Proposition}
\newtheorem{condition}{Condition}
\newtheorem{thm}{Theorem}
\newtheorem{defn}{Definition}
\newtheorem{cor}{Corollary}
\newtheorem{obs}{Observation}
 \numberwithin{equation}{section}
 

\setbeamertemplate{navigation symbols}{}


\title[PLSC 43401]{Mathematical Foundations of Political Methodology}

\author{Robert Gulotty}
\institute[Chicago]{University of Chicago}
\vspace{0.3in}




\begin{document}

\begin{frame}
\maketitle
\end{frame}





\begin{frame}
\frametitle{Topics for Today}
\begin{itemize}
\item LIE and LTV
\item Change of coordinates when Random Variables are Tranformed.
\end{itemize}
\end{frame}


\begin{frame}
\frametitle{Application of LIE: Conditional Expectation as an Estimator}
Suppose $Y$ is a random variable, and $X$ is some data.

\begin{eqnarray}
\hat{Y} & = & E[Y|X] \nonumber 
\end{eqnarray}\pause
So we know $\hat{Y}$ is an estimator.\\


\end{frame}


\begin{frame}
\frametitle{Estimation error}
Suppose $\hat{Y}$ is an estimator from before, the estimation error:
\begin{eqnarray}
\tilde{Y} & = & \hat{Y}-Y \nonumber 
\end{eqnarray}\pause
is a random variable where:
\begin{eqnarray}
E[\tilde{Y}|X] & = & E[ \hat{Y}-Y|X]=E[ \hat{Y}|X]-E[Y|X]=\hat{Y} -\hat{Y} =0\nonumber 
\end{eqnarray}\pause
By iterated expectations:\pause
\begin{eqnarray}
E[\tilde{Y}] & = & E[E[ \tilde{Y}|X]]=0 \nonumber 
\end{eqnarray}\pause
This estimator's error centers on 0.

\end{frame}


\begin{frame}
\frametitle{Correlation of  Estimator and Estimation Error}
Suppose $\hat{Y}$ is the estimator from before, and $\tilde{Y}$ the estimation error.
\begin{eqnarray}
E[\tilde{Y}\hat{Y}] & = & E[E[\tilde{Y}\hat{Y}|X]]=E[\hat{Y}E[\tilde{Y}|X]]=0 \nonumber 
\end{eqnarray}\pause
because $\hat{Y}$ only depends on $X$.\pause
\begin{eqnarray}
cov(\tilde{Y},\hat{Y}) & = & E[ \tilde{Y}\hat{Y}]-E[\tilde{Y}]E[\hat{Y}]=0- E[\tilde{Y}]0=0\nonumber 
\end{eqnarray}\pause
The error is uncorrelated with the estimate. 
\end{frame}


\begin{frame}
\frametitle{Law of Total Variance}
Suppose $X$ and $Y$ are random variables.  Then 

\begin{eqnarray}
var[X] & = & E[var[X|Y]]+var(E[X|Y]) \nonumber 
\end{eqnarray}
The first term is the average of the variation for each value of Y.\pause
The second term is the variation across the averages within each value of Y.
\end{frame}

\begin{frame}
\frametitle{Law of Total Variance: application}
Suppose $\theta$ is a parameter, and $x$ is data.

\begin{eqnarray}
var[\theta] & = & E[var[\theta|x]]+var(E[\theta|x]) \nonumber 
\end{eqnarray}
The expected variance of $\theta$ conditional on data (the posterior) is less than the variance in $\theta$ without the data (the prior).
\end{frame}


%
%
%\begin{frame}
%
%\begin{defn}
%$\boldsymbol{X} = (X_{1}, X_{2}, \hdots, X_{N}) $ is a vector of random variables.  If $\boldsymbol{X}$ has pdf 
%
%\begin{eqnarray}
%f(\boldsymbol{x}) & = & (2 \pi)^{-N/2} \text{det}\left(\boldsymbol{\Sigma}\right)^{-1/2} \exp\left(-\frac{1}{2}(\boldsymbol{x} - \boldsymbol{\mu})^{'}\boldsymbol{\Sigma} (\boldsymbol{x} - \boldsymbol{\mu} ) \right) \nonumber 
%\end{eqnarray}
%$\boldsymbol{X}$ is a \alert{Multivariate Normal} Distribution, 
%\begin{eqnarray}
%\boldsymbol{X} & \sim & \text{Multivariate Normal} (\boldsymbol{\mu}, \boldsymbol{\Sigma}) \nonumber 
%\end{eqnarray}
%
%
%\end{defn}
%
%
%\end{frame}
%
%
%\begin{frame}
%
%\begin{prop}
%Suppose $\boldsymbol{X} \sim \text{Multivariate Normal}(\boldsymbol{\mu}, \boldsymbol{\Sigma})$.  
%where $\boldsymbol{X}= (X_{1}, X_{2}, \hdots, X_{N})$. \\
%If $\text{cov}(X_{i}, X_{j})  = 0  $, then $X_{i}$ and $X_{j}$ are independent 
%
%\end{prop}
%\end{frame}
%
%
%
%\begin{frame}
%
%\begin{proof}
%Suppose $\mathbf{X}=(X_1,\ ... X_N)$ are MVN distributed \\ \pause 
%Partition these into two groups $\mathbf{X_A}$ and $\mathbf{X_B}$, with means $\boldsymbol{\mu_A}$ and $\boldsymbol{\mu_B}$, and variance covariance 
%$$\boldsymbol{\Sigma}=\ \begin{pmatrix} 
%\text{var}(X_{1}) & \text{cov}(X_{1}, X_{2}) & \hdots & \text{cov}(X_{1}, X_{N}) \\
%\text{cov}(X_{2}, X_{1}) & \text{var}(X_{2}) & \hdots & \text{cov}(X_{2}, X_{N} ) \\
%\vdots & \vdots & \ddots & \vdots \\
%\text{cov}(X_{N}, X_{1} ) & \text{cov}(X_{N}, X_{2} ) & \hdots & \text{var}(X_{N} ) \\
%\end{pmatrix}=\begin{pmatrix} \boldsymbol{\Sigma}_{AA} & \boldsymbol{\Sigma}_{AB} \\ \boldsymbol{\Sigma}_{BA} &\boldsymbol{\Sigma}_{BB} \\\end{pmatrix}$$
%
%\pause
%Now assume: $\boldsymbol{\Sigma}_{AB}=\boldsymbol{\Sigma}_{BA}=0$\pause
%\begin{eqnarray}
%\invisible<1-3>{& & [\mathbf{x_A^T}-\boldsymbol{\mu_A^T}, \quad \mathbf{x_B^T}-\boldsymbol{\mu_B^T}]\begin{bmatrix} \boldsymbol{\Sigma}_{AA} & 0\\ 0 &\boldsymbol{\Sigma}_{BB} \end{bmatrix}^{-1}  \begin{bmatrix}\mathbf{x_A}-\boldsymbol{\mu_A} \\ \mathbf{x_B}-\boldsymbol{\mu_B}\end{bmatrix} \nonumber } 
%\end{eqnarray}
%
%\end{proof}
%\end{frame}
%
%\begin{frame}
%Now taking that equation:
%\begin{eqnarray}
%\invisible<1>{& & [\mathbf{x_A^T}-\boldsymbol{\mu_A^T}, \quad \mathbf{x_B^T}-\boldsymbol{\mu_B^T}]\begin{bmatrix} \boldsymbol{\Sigma}_{AA} & 0\\ 0 &\boldsymbol{\Sigma}_{BB} \end{bmatrix}^{-1}  \begin{bmatrix}\mathbf{x_A}-\boldsymbol{\mu_A} \\ \mathbf{x_B}-\boldsymbol{\mu_B}\end{bmatrix} \nonumber \\ } \pause 
%\invisible<1-2>{& & [\mathbf{x_A^T}-\boldsymbol{\mu_A^T}, \quad \mathbf{x_B^T}-\boldsymbol{\mu_B^T}]\begin{bmatrix} \boldsymbol{\Sigma}_{AA}^{-1} & 0\\ 0 &\boldsymbol{\Sigma}_{BB}^{-1} \end{bmatrix}  \begin{bmatrix}\mathbf{x_A}-\boldsymbol{\mu_A} \\ \mathbf{x_B}-\boldsymbol{\mu_B}\end{bmatrix} \nonumber \\ } \pause 
%\invisible<1-3>{& &(\mathbf{x_A}-\boldsymbol{\mu_A})^T\boldsymbol{\Sigma}_{AA}^{-1}(\mathbf{x_A}-\boldsymbol{\mu_A})+ (\mathbf{x_B}-\boldsymbol{\mu_B})^{T} \boldsymbol{\Sigma}_{BB}^{-1}(\mathbf{x_B}-\boldsymbol{\mu_B}) \nonumber } \pause
%\end{eqnarray}
%
%\invisible<1-4>{So $exp\{(\mathbf{x}-\boldsymbol{\mu})^T\boldsymbol{\Sigma}^{-1}(\mathbf{x}-\boldsymbol{\mu})\}$ factorizes into :
%$$exp\{(\mathbf{x_A}-\boldsymbol{\mu_A})^T\boldsymbol{\Sigma}_{AA}^{-1}(\mathbf{x_A}-\boldsymbol{\mu_A})\}*exp\{(\mathbf{x_B}-\boldsymbol{\mu_B})^{T} \boldsymbol{\Sigma}_{BB}^{-1}(\mathbf{x_B}-\boldsymbol{\mu_B})\}$$}
%\pause
%so :\pause
%\invisible<1-5>{$$f(\mathbf{x})=g(\mathbf{x_A})h(\mathbf{x_B})$$}
%
%\end{frame}


\begin{frame}
\frametitle{Change of Coordinates}


\begin{prop}
Suppose $X$ is a random variable and $Y = g(X)$, where $g:\mathbb{R} \rightarrow \mathbb{R}$ that is a $monotonic$ function.  \\
Define $g^{-1}:\mathbb{R} \rightarrow \mathbb{R}$ such that $g^{-1}(g(X)) = X$ and is differentiable.  Then, 
\begin{eqnarray}
f_{Y}(y) & = & f_{X}(g^{-1}(y))\left|\frac{\partial g^{-1}(y)}{\partial y} \right| \text{ if $y = g(x)$ for some $x$ } \nonumber \\
& = & 0 \text{ otherwise } \nonumber 
\end{eqnarray}

\end{prop}


\end{frame}

\begin{frame}
\frametitle{Change of Coordinates}


\begin{proof}

Suppose $g(\cdot)$ is monotonically increasing (WLOG)\pause 
\begin{eqnarray}
F_{Y}(y) & = & P(Y \leq y) \nonumber \\ \pause 
& = & P(g(X) \leq y) \nonumber \\ \pause 
& = & P(X \leq g^{-1}(y) ) \nonumber \\ \pause 
& = & F_{X}(g^{-1}(y) ) \nonumber  \pause 
\end{eqnarray}
Now differentiating to get the pdf  \pause 
\begin{eqnarray}
\frac{\partial F_{Y}(y)}{\partial y} & = & \frac{\partial F_{X}(g^{-1}(y) )} {\partial y } \nonumber \\
 & = & f_{X}(g^{-1}(y) ) \frac{\partial g^{-1}(y)}{\partial y }   \nonumber  \pause 
\end{eqnarray}
Then this is a pdf because $\frac{\partial g^{-1}(Y)}{\partial y } > 0$.   


\end{proof}



\end{frame}


\begin{frame}
\frametitle{Application of Change of Coordinates}

Suppose $X$ is a random variable with pdf $f_{X}(x)$.  Suppose $Y =  X^{n}$.  Find $f_{Y}(y)$. \pause  \\
Then $g^{-1} (x)  =  x^{1/n}$.   \pause 
\begin{eqnarray}
f_{Y}(y) & = & f_{X}(g^{-1}(y)) \left| \frac{\partial g^{-1}(Y)}{\partial y } \right| \nonumber \\ \pause 
& = & f_{X} (y^{1/n} ) \frac{y^{\frac{1}{n} - 1}}{n} \nonumber 
\end{eqnarray}
\pause
\alert{We've used this to derive many of the pdfs} \pause 
\begin{itemize}
\item[-] General Normal distribution 
\item[-] Chi-Squared Distribution
\end{itemize}

\end{frame}

\begin{frame}
\frametitle{Example: Two Candidates}

Two candidates put out advertisements to persuade the median voter.  Suppose that each is off by a uniform error U(0,1).
\begin{itemize}
\item What is the PDF of the miss of the less persuasive advertisement?  What is its expectation?
\item Let X and Y be the error for the two candidates, and Z be the miss of the least persuasive advertisement.
$$Z=\max\{X,\ Y\}$$
\item $P(X\leq z)=P(Y\leq z) =z$
\item By independence of X and Y: $F_z(z)=P(\max\{X,\ Y\}\leqz)=P(X\leq z, Y \leq z)=P(X\leq z)P(Y \leq z)$
\item  $F_z(z)=z*z$
\item Take derivative, $f_z(z)=2z$
\item E(Z)=\int_0^1z 2z dz
\end{itemize}
\end{frame}



\begin{frame}
\frametitle{Getting to Central Limit Theorem}
\begin{itemize}
\item Use inequalities to prove Weak Law of Large Numbers (WLLN)
\item Use Taylor's Theorem to show how a Moment Generating Function (MGF) works.
\item  Use MGF and WLLN to prove Central Limit Theorem.
\end{itemize}

\end{frame}


\begin{frame}
\frametitle{Inequalities and Limit Theorems}

\alert{Limit Theorems}
\begin{itemize}
\item[-] What happens when we consider a long sequence of random variables?
\item[-] What can we reasonably infer from data?
\begin{itemize}
\item[-] Laws of large numbers: averages of random variables converge on expected value?
\item[-] Central Limit Theorems: sum of random variables have normal distribution?
\end{itemize}
\end{itemize}


\end{frame}


\begin{frame}
\frametitle{Weak Law of Large Numbers}

Proof plan: 
\begin{itemize}
\item[-] Markov's Inequality
\item[-] Chebyshev's Inequality
\item[-] Weak Law of Large Numbers 
\end{itemize}


\end{frame}



\begin{frame}
\frametitle{Markov's Inequality}

\begin{prop}
Suppose $X$ is a random variable that takes on non-negative values.  Then, for all $a>0$, 
\begin{eqnarray}
P(X\geq a) & \leq & \frac{E[X]}{a} \nonumber 
\end{eqnarray}
\end{prop}

\end{frame}


\begin{frame}
\frametitle{Markov's Inequality} 


\begin{proof}
For $a>0$,  \pause 
\begin{eqnarray}
E[X] & = & \int_{0}^{\infty} x f(x) dx \nonumber \\ \pause
& = & \int_{0}^{a} x f(x) dx + \int_{a}^{\infty} x f(x) dx \nonumber  \pause 
\end{eqnarray}

Because $X\geq 0$,  \pause 

\begin{eqnarray}
E[X] \geq \int_{a}^{\infty} x f(x) dx  \geq \int_{a}^{\infty} a f(x)dx = a P(X \geq a )\nonumber \\ \pause 
\frac{E[X]}{a} \geq P(X \geq a )\nonumber 
\end{eqnarray}


\end{proof}



\end{frame}

\begin{frame}
\frametitle{Chebyshev's Inequality} 

\begin{prop} 
If $X$ is a random variable with mean $\mu$ and variance $\sigma^2$, then, for any value $k>0$, 
\begin{eqnarray}
P(|X - \mu| \geq k) & \leq & \frac{\sigma^2}{k^2} \nonumber 
\end{eqnarray}

\end{prop}


\end{frame}

\begin{frame}
\frametitle{Chebyshev's Inequality} 


\begin{proof}
Define the random variable 
\begin{eqnarray}
 Y & = & (X - \mu)^2 \nonumber 
\end{eqnarray}

Where $\mu = E[X]$. 

\pause 

Then we know $Y$ is a non-negative random variable.  Set $a = k^2$.  \\ \pause 
Applying the inequality: \pause 
\begin{eqnarray}
P(Y \geq k^2) \leq \frac{E[Y]}{k^2} \nonumber \\ \pause 
P( (X- \mu)^2 \geq k^2) \leq \frac{E[(X - \mu)^2]}{k^2} \nonumber \\ \pause 
P( (X- \mu)^2 \geq k^2) \leq \frac{\sigma^2}{k^2}\nonumber 
\end{eqnarray}

\end{proof}

\end{frame}


\begin{frame}
\frametitle{Chebyshev's Inequality} 


Further we know that, 
\begin{eqnarray}
(X - \mu)^2  \geq k^2 \nonumber 
\end{eqnarray}

\pause 

Implies that  \pause 
\begin{eqnarray}
|X - \mu| \geq k \nonumber  \pause 
\end{eqnarray}

Thus, we have shown  \pause 
\begin{eqnarray}
P( |X- \mu| \geq k) \leq \frac{\sigma^2}{k^2}\nonumber 
\end{eqnarray}

\end{frame}



\begin{frame}{Example of Markov and Chebyshev}
Suppose that we know that the number of items produced in a factory is a random variable with mean 500.
\begin{itemize}
\item How likely are we to see a production of 1000?
\item If the variance is 50, how likely is the production between 400 and 600?
\end{itemize}

\end{frame}

\begin{frame}{Solution}
\begin{itemize}
\item How likely are we to see a production of 1000?
\item Markov Solution:
$$P(X\geq 1000)\leq\frac{E(X)}{1000}=\frac{500}{1000}$$\pause
\item If the variance is 50, how likely is the production between 400 and 600?
\item Chebyshev Inequality (set k to 100, $\sigma^2=50$):
$$P(|X-500| \geq 100)\leq\frac{\sigma^2}{(100)^2}=\frac{50}{10000}=\frac{1}{200}$$
$$P(|X-500| < 100)\geq 1-\frac{1}{200}$$
\end{itemize}

\end{frame}

\begin{frame}
\frametitle{Random Variables, Data and the True parameters}
\begin{itemize}
\item Capital letters X, Y, Z are random variables.
\item Greek letters $\mu,\ \sigma, \alpha, \beta, \lambda, \theta, \delta$ are the True parameters.\pause
\item Observed facts about the world are not the TRUTH - we see only shadows!
\begin{itemize}\pause
\item We infer the true value by observing data, but we can only approximate the truth. 
\item e.g. Fundamental problem causal inference $\delta=Y_t(i)-Y_c(i)$
\end{itemize}
\item To recognize this, we model our data as random variables, summaries of which are also random variables.
$$\bar{X}=\frac{X_{1}+ X_{2}+ \hdots+ X_{n}}{n}$$
\item $\bar{X}\neq \mu$
\end{itemize}

\end{frame}


\begin{frame}
\frametitle{Sequences of Random Variables}

Sequence of Independent and Identically Distributed Random variables.
\begin{itemize}
\item[-] Sequence: $X_{1}, X_{2}, \hdots, X_{n}, \hdots $
\item[-] Think of a sequence as sampled \alert{data}:
\begin{itemize}
\item[-] Suppose we are drawing a sample of $N$ observations
\item[-] Each observation will be a random variable, say $X_{i}$
\item[-] With realization $x_{i}$ 
\end{itemize}
\end{itemize}

\end{frame}


\begin{frame}
\frametitle{Mean/Variance of Sample Mean}

\begin{prop}
Let $X_{1}, X_{2}, \hdots, X_{n}$ be a random sample from a distribution with mean $\mu$ and variance $\sigma^2$.  Let $\bar{X}_{n}$ be the sample mean.  Then 
$E[\bar{X}_{n}] = \mu$ and var$(\bar{X}_{n}) = \frac{\sigma^2}{n}$
\end{prop}

\pause 
\begin{proof}

\begin{eqnarray}
E[\bar{X}_{n}] & = & \frac{1}{n}\sum_{i=1}^{n} E[X_{i}] \nonumber \\ \pause 
& = & \frac{1}{n} n \mu  = \mu \nonumber 
\end{eqnarray}



\end{proof}


\end{frame}


\begin{frame}
\frametitle{Mean/Variance of Sample Mean}


\pause 
\begin{eqnarray}
\text{var}(\bar{X}_{n}) & = & \frac{1}{n^2} \text{var}(\sum_{i=1}^{n} X_{i}) \nonumber \\ \pause 
 &= & \frac{1}{n^2} \sum_{i=1}^{n} \text{var}(X_{i}) \nonumber \\ \pause 
 & = & \frac{1}{n^2} n \sigma^2 = \frac{\sigma^2}{n} \nonumber  
\end{eqnarray}


\end{frame}



\begin{frame}
\frametitle{Weak Law of Large Numbers} 

\begin{prop}
Suppose $X_{1}, X_{2}, \hdots, X_{n}$ is a random sample from a distribution with mean $\mu$ and $Var(X_{i})= \sigma^2$.  Then, for all $\epsilon >0$, 
\begin{eqnarray}
P\left\{ \left| \frac{X_{1} + X_{2} + \hdots + X_{n} }{n} -\mu \right| \geq \epsilon \right\} \rightarrow 0 \text{ as } n \rightarrow \infty \nonumber 
\end{eqnarray}

\end{prop}


\end{frame}


\begin{frame}
\frametitle{Weak Law of Large Numbers}

\begin{proof}
From our previous proposition \pause 
\begin{eqnarray}
\invisible<1>{\frac{\text{E}[X_{1} + X_{2} + \cdots + X_{n} ]}{n} & = & \frac{\sum_{i=1}^{n} E[X_{i}] }{n}  = \mu \nonumber} \pause 
\end{eqnarray}
\invisible<1-2>{Further, } \pause 
\begin{eqnarray}
\invisible<1-3>{\text{E}[ (\frac{\sum_{i=1}^{n} X_{i} - \mu}{n} )^2] = \frac{\text{Var}(X_{1} + X_{2} + \cdots + X_{n} )}{n^2} & = & \frac{ \sum_{i=1}^{n} \text{Var}(X_{i}) }{n^2} = \frac{\sigma^2}{n} \nonumber } \pause 
\end{eqnarray}
\invisible<1-4>{Apply Chebyshev's Inequality: } \pause 
\begin{eqnarray}
\invisible<1-5>{P\left\{ \left| \frac{X_{1} + X_{2} + \hdots + X_{n} }{n} -\mu \right| \geq \epsilon \right\}\leq \frac{\sigma^2}{n \epsilon^2} \nonumber }
\end{eqnarray}

\end{proof}


\end{frame}

\begin{frame}
Suppose $X_{1}, X_{2}, \hdots$ are iid normal distributions,
\begin{eqnarray}
X_{i} & \sim & \text{Normal}(0, 10) \nonumber 
\end{eqnarray}

\begin{eqnarray} 
P\left\{ \left| \frac{X_{1} + X_{2} + \hdots + X_{n} }{n} -\mu \right| \geq 0.1 \right\} \text{ as } n \rightarrow \infty\nonumber 
\end{eqnarray}
\pause 

Suppose we want to guarantee that we have at most a $0.01$ probability of being more than $0.1$ away from the true $\mu$.  How big do we need $n$?\pause 
\begin{eqnarray}
0.01 & = & \frac{10}{n (0.1^2) } \nonumber \pause  \\
n & = & \frac{1000}{0.01} \nonumber  \pause \\
n & = & 100,000  \nonumber 
\end{eqnarray}



\end{frame}

\end{document}